\slajdid{10-s086}
\begin{frame}[label=10-s086]
\gusttekst\setlength{\abovedisplayskip}{3pt}\setlength{\belowdisplayskip}{3pt}\setlength{\abovedisplayshortskip}{2pt}\setlength{\belowdisplayshortskip}{2pt}\begin{columns}[T,onlytextwidth]\column{.27\textwidth}\centering\input{slike/tikz/s082_ecl_diferencijalni_par.tex}\column{.33\textwidth}\centering\begin{tikzpicture}[x=.65cm,y=.8cm,font=\sitneoznake,>=Latex]
\draw[->](0,0)--(5,0)node[below]{$V_I$};\draw[->](.2,-.1)--(.2,2.8)node[above]{$V_{O1},V_{O2}$};
\draw[dashed](.2,2.1)--(1.8,2.1);\draw[dashed](.2,.55)--(1.8,.55);\node[left]at(.2,2.1){$V_{OH}$};\node[left]at(.2,.55){$V_{OL}$};
\draw(.65,2.1)--(1.8,2.1)--(2.7,.55)--(3.8,.55)--(4.5,1.25);\draw(.65,.55)--(1.8,.55)--(2.7,2.1)--(4.5,2.1);\node[above]at(3.5,2.1){$V_{O2}$};\node[above]at(3.6,.55){$V_{O1}$};
\foreach\x/\l/\y in{1.8/V_{IL}/2.1,2.25/V_R/1.325,2.7/V_{IH}/.55}{\draw[dashed](\x,0)--(\x,\y);}
\node[below left]at(1.8,0){$V_{IL}$};\node[below]at(2.25,-.4){$V_R$};\draw(2.25,0)--(2.25,-.35);\node[below right]at(2.7,0){$V_{IH}$};\draw[->,DEplava](4.85,1.65)--(4.22,1.12);\end{tikzpicture}

\[V_{OL}=V_{CC}-R_C I_E\]
\[V_{OH}=V_{CC}\]
\[V_{IL}=V_{\gamma T}-V_{BE}+V_R\]
\[V_{IH}=V_{BE}-V_{\gamma T}+V_R\]
\column{.38\textwidth}
Ovo je neregularan rad ECL kola. Kao što smo videli do ovog ulaznog napona tranzistori T1 i T2 su ili zakočeni ili rade u aktivnom režimu. Time je obezbeđeno i da mogu brzo da se uključe ali da mogu i da se brzo isključe. Time je dobijena velika brzina rada ECL logičkih kola. Znači ovu situaciju da tranzistor T1 ide u zasićenje treba izbegavati.
\end{columns}
\end{frame}
